\chapter*{\uppercase{ABSTRACT}}
\addcontentsline{toc}{section}{\bfseries \uppercase{Abstract}}
AI agents that don't have persistent memory forget instructions and project
decisions once a session ends, and when an instruction changes they often keep using
the old one. For this project we built a long-term memory engine in Rust for such agents.
What it remembers carries over between sessions, can be corrected later, and can be
traced back to the original messages.

The engine has five modules: source acquisition and processing, knowledge
extraction and consolidation, temporal knowledge and graph management, hybrid
retrieval and grounded reasoning, and model training and evaluation. Original messages
are stored unchanged as source events. Assertions (typed statements about the user or
project) are extracted along with the exact source quote they came from and stored as
versions, so a correction replaces an old value without deleting it.
Retrieval combines keyword, vector, time and graph search using Reciprocal Rank
Fusion and returns each result with its source quotes. The extraction and answering is done by a small local
Qwen3-4B model that we fine-tuned with QLoRA.

We tested the engine on LoCoMo, a public benchmark of ten long conversations (1,540 questions), using the
same reader model for every setup and two model judges. With the engine the reader answered 83.1\% of the
questions correctly, compared with 80.0\% for plain retrieval over the raw text, 43.0\% when keeping only the
last 8,192 tokens, and 0.3\% with no memory. Giving the reader the whole conversation scored 89.0\%, but used more
than three times as many tokens. All 99,971 source quotes returned by the engine matched the conversation
exactly. Our own scripted tests showed that corrections keep their history and that the engine refuses to answer
without evidence, and they also caught a date bug, which we fixed after the test.
